\documentclass[fontsize=12pt,
               paper=a4,
               twoside=false,
               parskip=half,
               ]{scrartcl}

% Load the packages
%   Copyright 2012 Comet Engineering, Patrick Haring & Christian Bürgi
%
%   Licensed under the Apache License, Version 2.0 (the "License");
%   you may not use this file except in compliance with the License.
%   You may obtain a copy of the License at
%
%       http://www.apache.org/licenses/LICENSE-2.0
%
%   Unless required by applicable law or agreed to in writing, software
%   distributed under the License is distributed on an "AS IS" BASIS,
%   WITHOUT WARRANTIES OR CONDITIONS OF ANY KIND, either express or implied.
%   See the License for the specific language governing permissions and
%   limitations under the License.

% Packages Template
% =================
% 
% Contains packages used for project documentation
% 
% @author burgc5
% 
% To use this simply enter: %   Copyright 2012 Comet Engineering, Patrick Haring & Christian Bürgi
%
%   Licensed under the Apache License, Version 2.0 (the "License");
%   you may not use this file except in compliance with the License.
%   You may obtain a copy of the License at
%
%       http://www.apache.org/licenses/LICENSE-2.0
%
%   Unless required by applicable law or agreed to in writing, software
%   distributed under the License is distributed on an "AS IS" BASIS,
%   WITHOUT WARRANTIES OR CONDITIONS OF ANY KIND, either express or implied.
%   See the License for the specific language governing permissions and
%   limitations under the License.

% Packages Template
% =================
% 
% Contains packages used for project documentation
% 
% @author burgc5
% 
% To use this simply enter: %   Copyright 2012 Comet Engineering, Patrick Haring & Christian Bürgi
%
%   Licensed under the Apache License, Version 2.0 (the "License");
%   you may not use this file except in compliance with the License.
%   You may obtain a copy of the License at
%
%       http://www.apache.org/licenses/LICENSE-2.0
%
%   Unless required by applicable law or agreed to in writing, software
%   distributed under the License is distributed on an "AS IS" BASIS,
%   WITHOUT WARRANTIES OR CONDITIONS OF ANY KIND, either express or implied.
%   See the License for the specific language governing permissions and
%   limitations under the License.

% Packages Template
% =================
% 
% Contains packages used for project documentation
% 
% @author burgc5
% 
% To use this simply enter: \input{./packages.tex}

\usepackage[utf8]{inputenc}
\usepackage[T1]{fontenc}

% Set font to latin modern
\usepackage{lmodern}

\usepackage[pdftex]{graphicx}
\usepackage{epstopdf}

% Create links in pdf documents
\usepackage[colorlinks,pdfpagelabels,pdfstartview=FitH,bookmarksopen=true,bookmarksnumbered=true,linkcolor=black,plainpages=false,hypertexnames=false,citecolor=black] {hyperref}
\hypersetup{
    colorlinks,%
    citecolor=black,%
    filecolor=black,%
    linkcolor=black,%
    urlcolor=black
}
\urlstyle{same}

% Use \enquote{} to create quotation marks
\usepackage{csquotes}

% Create professional tables with booktabs
% @see http://en.wikibooks.org/wiki/LaTeX/Tables#Professional_tables
\usepackage{booktabs}

% Customizable enumerates/itemizes
\usepackage{enumitem}

% git meta information
\usepackage{gitinfo}


\usepackage[utf8]{inputenc}
\usepackage[T1]{fontenc}

% Set font to latin modern
\usepackage{lmodern}

\usepackage[pdftex]{graphicx}
\usepackage{epstopdf}

% Create links in pdf documents
\usepackage[colorlinks,pdfpagelabels,pdfstartview=FitH,bookmarksopen=true,bookmarksnumbered=true,linkcolor=black,plainpages=false,hypertexnames=false,citecolor=black] {hyperref}
\hypersetup{
    colorlinks,%
    citecolor=black,%
    filecolor=black,%
    linkcolor=black,%
    urlcolor=black
}
\urlstyle{same}

% Use \enquote{} to create quotation marks
\usepackage{csquotes}

% Create professional tables with booktabs
% @see http://en.wikibooks.org/wiki/LaTeX/Tables#Professional_tables
\usepackage{booktabs}

% Customizable enumerates/itemizes
\usepackage{enumitem}

% git meta information
\usepackage{gitinfo}


\usepackage[utf8]{inputenc}
\usepackage[T1]{fontenc}

% Set font to latin modern
\usepackage{lmodern}

\usepackage[pdftex]{graphicx}
\usepackage{epstopdf}

% Create links in pdf documents
\usepackage[colorlinks,pdfpagelabels,pdfstartview=FitH,bookmarksopen=true,bookmarksnumbered=true,linkcolor=black,plainpages=false,hypertexnames=false,citecolor=black] {hyperref}
\hypersetup{
    colorlinks,%
    citecolor=black,%
    filecolor=black,%
    linkcolor=black,%
    urlcolor=black
}
\urlstyle{same}

% Use \enquote{} to create quotation marks
\usepackage{csquotes}

% Create professional tables with booktabs
% @see http://en.wikibooks.org/wiki/LaTeX/Tables#Professional_tables
\usepackage{booktabs}

% Customizable enumerates/itemizes
\usepackage{enumitem}

% git meta information
\usepackage{gitinfo}


% SVN Meta
\SVN $Date: 2012-05-11 14:51:38 +0200 (Fr, 11 Mai 2012) $
\SVN $Revision: 101 $
\SVN $HeadURL: https://svn.bfh.ch/repos/projects/patmon1/trunk/doc/src/use_case_model.tex $


\begin{document}

% Document title for title.tex
\newcommand{\doctitle}{Use Case Model}
% Titlepage Template
% ==================
% 
% @author burgc5
% 
% To use this simply enter: % Titlepage Template
% ==================
% 
% @author burgc5
% 
% To use this simply enter: % Titlepage Template
% ==================
% 
% @author burgc5
% 
% To use this simply enter: \input{./title.tex}
% 
% You have to define the commands '\doctitle' and '\docrevision' to give the 
% document a title and a revision on its titlepage.
% Do this with the following command:
% \newcommand{\doctitle}{Document title goes here}
%
% SVN:
% ----
% You also have to define the variables:
% \SVN $Date$
% \SVN $Revision$
%
% As executing the following command on the file:
% > svn propset svn:keywords "Date Revision" filename.tex
% 
% This titlepage needs:
% \usepackage[pdftex]{graphicx}
% \usepackage{svn}
%

\begin{titlepage}

\begin{center}

% Team-logo
\includegraphics[width=0.35\textwidth]{./comet-logo.eps}\\[2.5cm]    

% Project title
\textsc{\Large Comet Pinball}\\[2cm]

% Document title
{ \huge \bfseries \doctitle{}}\\[3cm]

% Members/Client
\begin{minipage}{0.45\textwidth}
\begin{flushleft} \large
\emph{Team Members:}\\
Patrick \textsc{Haring}\\
Christian \textsc{Bürgi}
\end{flushleft}
\end{minipage}
\begin{minipage}{0.45\textwidth}
\begin{flushright} \large
\emph{Client:} \\
Jean-Pierre \textsc{Caillot}\\
~
\end{flushright}
\end{minipage}

\vfill

{\large 
Revision hash: \gitAbbrevHash \\[0.2cm]
Commit time: \gitCommitterIsoDate \\[0.2cm]
{\footnotesize \itshape \url{https://github.com/boskoop/comet-pinball/}}}

\end{center}

\end{titlepage}
% 
% You have to define the commands '\doctitle' and '\docrevision' to give the 
% document a title and a revision on its titlepage.
% Do this with the following command:
% \newcommand{\doctitle}{Document title goes here}
%
% SVN:
% ----
% You also have to define the variables:
% \SVN $Date$
% \SVN $Revision$
%
% As executing the following command on the file:
% > svn propset svn:keywords "Date Revision" filename.tex
% 
% This titlepage needs:
% \usepackage[pdftex]{graphicx}
% \usepackage{svn}
%

\begin{titlepage}

\begin{center}

% Team-logo
\includegraphics[width=0.35\textwidth]{./comet-logo.eps}\\[2.5cm]    

% Project title
\textsc{\Large Comet Pinball}\\[2cm]

% Document title
{ \huge \bfseries \doctitle{}}\\[3cm]

% Members/Client
\begin{minipage}{0.45\textwidth}
\begin{flushleft} \large
\emph{Team Members:}\\
Patrick \textsc{Haring}\\
Christian \textsc{Bürgi}
\end{flushleft}
\end{minipage}
\begin{minipage}{0.45\textwidth}
\begin{flushright} \large
\emph{Client:} \\
Jean-Pierre \textsc{Caillot}\\
~
\end{flushright}
\end{minipage}

\vfill

{\large 
Revision hash: \gitAbbrevHash \\[0.2cm]
Commit time: \gitCommitterIsoDate \\[0.2cm]
{\footnotesize \itshape \url{https://github.com/boskoop/comet-pinball/}}}

\end{center}

\end{titlepage}
% 
% You have to define the commands '\doctitle' and '\docrevision' to give the 
% document a title and a revision on its titlepage.
% Do this with the following command:
% \newcommand{\doctitle}{Document title goes here}
%
% SVN:
% ----
% You also have to define the variables:
% \SVN $Date$
% \SVN $Revision$
%
% As executing the following command on the file:
% > svn propset svn:keywords "Date Revision" filename.tex
% 
% This titlepage needs:
% \usepackage[pdftex]{graphicx}
% \usepackage{svn}
%

\begin{titlepage}

\begin{center}

% Team-logo
\includegraphics[width=0.35\textwidth]{./comet-logo.eps}\\[2.5cm]    

% Project title
\textsc{\Large Comet Pinball}\\[2cm]

% Document title
{ \huge \bfseries \doctitle{}}\\[3cm]

% Members/Client
\begin{minipage}{0.45\textwidth}
\begin{flushleft} \large
\emph{Team Members:}\\
Patrick \textsc{Haring}\\
Christian \textsc{Bürgi}
\end{flushleft}
\end{minipage}
\begin{minipage}{0.45\textwidth}
\begin{flushright} \large
\emph{Client:} \\
Jean-Pierre \textsc{Caillot}\\
~
\end{flushright}
\end{minipage}

\vfill

{\large 
Revision hash: \gitAbbrevHash \\[0.2cm]
Commit time: \gitCommitterIsoDate \\[0.2cm]
{\footnotesize \itshape \url{https://github.com/boskoop/comet-pinball/}}}

\end{center}

\end{titlepage}

\tableofcontents

%=======
% ACTORS
%=======

\section{Actors}

\subsection{List of actors}

\begin{tabular}{llp{9cm}}
\toprule
\textbf{Actor} & \textbf{Type} & \textbf{Description} \\ 
\midrule
\emph{Player} & primary & \emph{Player}s run the Pinball Simulator on their computer.\\
%\emph{Administrator} & primary & The \emph{administrator} is an employee of the hospital as
%well. He is responsible for the hospital's IT systems to work. He configures
%the system for the \emph{doctor}s to work with. He is not a \emph{doctor}. \\ 
%\emph{OT-Logger} & supporting & The \emph{OT-Logger}s are devices which are used to monitor
%\emph{patient}'s temperatures over a period of time. They can be remotely polled by 
%the system. A device can only be reused after appliance on a \emph{patient} if the device was reset.\\ 
%\emph{Patient} & off-stage & \emph{Patient}s have the monitoring device attached for a period from a few days to a few weeks. \\
%\emph{Email-Server} & supporting & The \emph{email-server} will send the activation code for the first authentication of the \emph{doctor} against the server.

\\ 
\bottomrule 
\end{tabular} 

\subsection{Primary actor goals}

\begin{description}

\item[Player:] \hfill \\
A \emph{Player} uses the Pinball Simulator to pass his leisure time and enjoy himself. He is interested in a well designed play field with nice visual and acoustical effects.

%\item[Administrator:] \hfill \\
%The \emph{administrator} wants to create and delete accounts for the \emph{doctor}s.

\end{description}

%==================
% Use cases (brief)
%==================


\section{Use cases (brief)}

\subsection{Player}

\begin{description}

\item[Start game:]  The application has to be running on a players computer. The player wants to start a game so he enters the main screen of the application and 


\item[Ball leaves via drain]

\item[Extra ball:]


\item[Extra game:]

\item[Gameover:]

\item[Ball is stuck:]

\item[Obstacle on the play field is hit:]

\item[Tilt:]

\item[:]

\end{description}







\subsection{Administrator}

\begin{description}	
\item[Create account:] A new \emph{doctor} has to be registered in the system. The \emph{administrator} enters the \emph{doctor}'s name and mail address. The system presents the entered information. The \emph{administrator} verifies it and confirms that to the system. The system creates an account with a random initial password and presents after successful creation the full account data to the \emph{administrator}. He records the initial password which he has to give to the \emph{doctor} in person. There is also an email sent to the entered account email by using the \emph{email-server}. It contains an activation code, which is used by the \emph{doctor} for his first login.

%\item[Delete account:] A \emph{doctor} has to be deleted from the system. The account can only be deleted if it was disabled (and not enabled again) for a given period. The \emph{administrator} archives the data of the \emph{patient}s, which are referenced to this account and deletes the \emph{doctor}'s account. By deleting a \emph{doctor}'s account the system also deletes all corresponding \emph{patient} data (in the system).

%\item[Disable account:] An account of a \emph{doctor} has to be disabled. The \emph{administrator} disables the login of the \emph{doctor}. Thereby the account is also deactivated (when enabled again, the \emph{doctor} needs to reactivate the account). The system notifies the \emph{doctor} by means of an email via the \emph{email-server}. 

%\item[Enable account:] A disabled account of a \emph{doctor} has to be enabled again. The System generates a new activation code and sends it via the \emph{email-server} to the \emph{doctor}. If necessary, he also resets the \emph{doctor}'s password.

%\item[Reset password:] A \emph{doctor}'s password has to be changed. The \emph{administrator} generates a new password and gives it to the \emph{doctor} in person.

%\item[Authenticate against System:] The \emph{administrator} authenticates against the system by entering his email address and password.

\end{description}

%==========================
% Use cases (fully-dressed)
%==========================

\section{Use cases (fully-dressed)}

% UC1
%----

\subsection{Use Case UC1: Assign device}

\textbf{\textsf{Scope:}} Comet Patient Monitor

%\textbf{\textsf{Level:}} user goal

%\textbf{\textsf{Primary actor:}} \emph{Doctor}

%\textbf{\textsf{Preconditions:}} \emph{Doctor} is authenticated against the system. \emph{Patient} is registered in the System.

%\textbf{\textsf{Postconditions:}} The monitoring device is assigned to the \emph{patient}. The device will monitor the temperature during a given period.

%\textbf{\textsf{Main success scenario:}}

%\begin{enumerate}[leftmargin=3em]
%	\item \emph{Doctor} enters the \emph{patient}'s name and/or social security number into system.
%	\item System looks up matching \emph{patient} profiles and displays them to the \emph{doctor}.
%	\item \emph{Doctor} selects the appropriate \emph{patient} profile.
%	\item System displays detailed version of the selected profile.
%	\item \emph{Doctor} enters the device identification number into the system.
%	\item System checks if the device is already assigned to a \emph{patient}.
%	\item \emph{Doctor} enters the beginning and the end of the observation period into the system.
%	\item System verifies that the entered data is a valid observation period.
%	\item \emph{Doctor} enters the frequency by which the measures have to be performed.
%	\item System configures the device.
%	\item System saves assignment of the device to the \emph{patient}.
%\end{enumerate}

%\textbf{\textsf{Extensions:}}

\begin{itemize}[leftmargin=3em]
	\item[*a.] At any time, \emph{doctor} cancels assignment operation.
	\begin{enumerate}
		\item System discards all entered data.
	\end{enumerate}
	\item[2a.] No profile is found.
	\begin{enumerate}
		\item System signals error and reverts to step 1 in main scenario.
	\end{enumerate}
	\item[2b.] Correct \emph{patient} profile is not found.
	\begin{enumerate}
		\item \emph{Doctor} verifies entered data.
		\begin{itemize}
			\item[1a.] \emph{Patient} is not registered (violation of precondition!):
			\begin{enumerate}[label=\arabic*.]
				\item \emph{Doctor} cancels assignment operation and proceeds with UC4.
			\end{enumerate}
		\end{itemize}
		\item \emph{Doctor} proceeds with step 1 in main scenario.
	\end{enumerate}
	\item[4a.] \emph{Doctor} selected wrong \emph{patient} profile.
	\begin{enumerate}
		\item \emph{Doctor} cancels assignment to selected profile.
		\item System reverts to step 2 in main scenario.
	\end{enumerate}
	\item[6a.] Device is already assigned.
	\begin{enumerate}
		\item System signals error.
		\item \emph{Doctor} gets a different device.
		\item System reverts to step 5 in main scenario.
	\end{enumerate}
	\item[8a.] Invalid observation period.
	\begin{enumerate}
		\item System signals error and reverts to step 7 in main scenario.
	\end{enumerate}
	\item[10a.] Error while configuring the device.
	\begin{enumerate}
		\item System signals error.
		\item \emph{Doctor} marks device as faulty and sends it to reparation.
		\item \emph{Doctor} gets a different device.
		\item System reverts to step 5 in main scenario.
	\end{enumerate}
\end{itemize}

\section{Use case diagram}


% Diagram generated with astah
%\includegraphics[width=15cm]{./use-case-diagram.png}

\end{document}